\section{Stakeholders}

Sono considerati tutti i soggetti che hanno un interesse diretto o indiretto
nel sistema. Per ognuno è indicato il ruolo e l'interesse specifico.

\begin{table}[H]
    \centering
    \begin{tabularx}{\textwidth}{@{}lX@{}}
        \toprule
        \textbf{Stakeholder} & \textbf{Descrizione e interesse} \\
        \midrule
        Utente anonimo &
        Visitatore non registrato che consulta il catalogo, cerca brani per
        titolo o artista e ne legge il testo (piano o sincronizzato).
        Può contribuire pubblicando nuovi testi superando la sfida PoW. \\
        \addlinespace
        Utente registrato &
        Possessore di un account autenticato via JWT. Pubblica testi senza
        risolvere la PoW e i contributi gli vengono attribuiti
        (\emph{submittedBy}). Può modificare o cancellare i propri brani. \\
        \addlinespace
        Amministratore di sistema &
        Responsabile dell'esercizio del backend: deployment, backup,
        monitoraggio dei log, gestione delle chiavi crittografiche
        (\code{JWT\_SECRET\_KEY}, \code{SECRET\_KEY}). \\
        \addlinespace
        Sviluppatore &
        Mantiene il codice di backend e mobile app; ha interesse a una base
        di codice modulare, testabile e ben documentata. \\
        \addlinespace
        Autori e detentori dei diritti &
        Titolari dei diritti d'autore sui testi pubblicati. Hanno interesse a
        che il sistema rispetti la normativa sul copyright e supporti
        procedure di rimozione (takedown) dei contenuti. \\
        \addlinespace
        Garante per la protezione dei dati personali &
        Vigila sul rispetto del GDPR per i dati degli utenti registrati
        (username, email, hash della password). \\
        \bottomrule
    \end{tabularx}
    \caption{Stakeholders del sistema Stopify.}
\end{table}
